\documentclass{article}
\usepackage{amsmath,amssymb,geometry}
\geometry{a4paper,margin=1.5cm}
\usepackage{ctex}
\usepackage{enumitem}
\usepackage{xcolor}

\newcommand{\comment}[1]{\textcolor{blue}{\small\textbf{【指导】#1}}}

\begin{document}
	
	\begin{enumerate}[label=\textbf{(\arabic*)},leftmargin=*]
	\item
	\begin{enumerate}[label=\textbf{\arabic*.},leftmargin=*]
		\item \textbf{建系与坐标化}：以 $D$ 为原点，$DA, DC, DD_1$ 所在直线分别为 $x, y, z$ 轴建立空间直角坐标系。\\
		各点坐标：
		\[
		\begin{array}{c}
			D(0,0,0),\; A(2,0,0),\; B(2,2,0),\; C(0,2,0),\\
			A_1(2,0,2),\; B_1(2,2,2),\; C_1(0,2,2),\; D_1(0,0,2)
		\end{array}
		\]
		\comment{书写坐标时注意 $xOy$ 平面内的点竖坐标为 $0$，上底面各点竖坐标均为 $2$。}
		
		\item \textbf{证明线线平行}：
		\[
		\overrightarrow{A_1D} = (0-2,0-0,0-2)=(-2,0,-2),\quad 
		\overrightarrow{B_1C} = (0-2,2-2,0-2)=(-2,0,-2)
		\]
		观察得 $\overrightarrow{A_1D} = -\overrightarrow{BC_1}$，故 $\overrightarrow{A_1D} \parallel \overrightarrow{B_1C}$，即直线 $A_1D \parallel B_1C$。\\
		\comment{向量共线是证明线线平行的核心依据，注意方向向量的选取。}
		
		\item \textbf{证明线面垂直}：
		\[
		\overrightarrow{BD_1} = (0-2,0-2,2-0)=(-2,-2,2)
		\]
		取平面 $A_1C_1D$ 内的两条相交直线 $A_1C_1$ 和 $A_1D$：
		\[
		\overrightarrow{A_1C_1} = (0-2,2-0,2-2)=(-2,2,0),\quad 
		\overrightarrow{A_1D} = (-2,0,-2)
		\]
		计算数量积：
		\[
		\begin{aligned}
			\overrightarrow{BD_1} \cdot \overrightarrow{A_1C_1} &= (-2)\times(-2)+(-2)\times2+2\times0 = 4-4+0=0,\\
			\overrightarrow{BD_1} \cdot \overrightarrow{A_1D} &= (-2)\times(-2)+(-2)\times0+2\times(-2) = 4+0-4=0.
		\end{aligned}
		\]
		因此 $BD_1 \perp A_1C_1$，$BD_1 \perp A_1D$，且 $A_1C_1 \cap A_1D = A_1$，故 $BD_1 \perp$ 平面 $A_1C_1D$。\\
		\comment{线面垂直需证直线垂直于平面内两条相交直线，数量积为 $0$ 是垂直的代数表达。}
	\end{enumerate}
	
	
	\item
	\begin{enumerate}[label=\textbf{\arabic*.},leftmargin=*]
		\item \textbf{坐标化}：沿用原坐标系，各点坐标：
		\[
		\begin{array}{c}
			A(2,0,0),\; B(2,2,0),\; C(0,2,0),\; D(0,0,0),\\
			A_1(2,0,2),\; B_1(2,2,2),\; C_1(0,2,2),\; D_1(0,0,2),\; E(0,0,1)
		\end{array}
		\]
		
		\item \textbf{求平面 $AEC_1$ 的法向量}：
		\[
		\overrightarrow{AE} = (0-2,0-0,1-0)=(-2,0,1),\quad 
		\overrightarrow{AC_1} = (0-2,2-0,2-0)=(-2,2,2)
		\]
		设 $\boldsymbol{n}=(x,y,z)$，由 $\boldsymbol{n}\perp\overrightarrow{AE}$ 且 $\boldsymbol{n}\perp\overrightarrow{AC_1}$：
		\[
		\begin{cases}
			-2x + 0\cdot y + 1\cdot z = 0 \quad \Rightarrow\quad z = 2x\\
			-2x + 2y + 2z = 0
		\end{cases}
		\]
		代入 $z=2x$ 到第二式：$-2x+2y+4x=0 \Rightarrow 2x+2y=0 \Rightarrow y = -x$。
		取 $x=1$，得 $y=-1,\;z=2$，故 $\boldsymbol{n}=(1,-1,2)$。
		\comment{法向量的坐标通常取整数或简单分数，便于后续计算。}
		
		\item \textbf{线面角的正弦值}：
		$\overrightarrow{A_1B} = (2-2,2-0,0-2) = (0,2,-2)$。
		设直线 $A_1B$ 与平面 $AEC_1$ 所成角为 $\theta$，则
		\[
		\sin\theta = |\cos\langle\overrightarrow{A_1B},\boldsymbol{n}\rangle|
		= \frac{|\overrightarrow{A_1B}\cdot\boldsymbol{n}|}{|\overrightarrow{A_1B}|\cdot|\boldsymbol{n}|}
		= \frac{|0\times1 + 2\times(-1) + (-2)\times2|}{\sqrt{0^2+2^2+(-2)^2}\cdot\sqrt{1^2+(-1)^2+2^2}}
		= \frac{|-2-4|}{\sqrt{8}\cdot\sqrt{6}} = \frac{6}{\sqrt{48}} = \frac{6}{4\sqrt{3}} = \frac{\sqrt{3}}{2}.
		\]
		故 $\sin\theta = \dfrac{\sqrt{3}}{2}$，即 $\theta = 60^\circ$。
		\comment{线面角的正弦值等于方向向量与法向量夹角余弦的绝对值。}
		
		\item \textbf{二面角的余弦值}：
		平面 $ABCD$ 的法向量 $\boldsymbol{m}=(0,0,1)$。
		\[
		\cos\langle\boldsymbol{m},\boldsymbol{n}\rangle = \frac{(0,0,1)\cdot(1,-1,2)}{1\cdot\sqrt{6}} = \frac{2}{\sqrt{6}} = \frac{\sqrt{6}}{3}.
		\]
		观察图形，平面 $AEC_1$ 与底面 $ABCD$ 所成二面角为锐角，故二面角的余弦值为 $\dfrac{\sqrt{6}}{3}$。
		
		\comment{二面角需结合图形判断锐钝，不能直接用法向量夹角。}
	\end{enumerate}
	
	
	\item 详见作业提升题
	
	\end{enumerate}
	
\end{document}